% Full-model signed certificates. No arithmetic classification of mu is asserted.
% Dependencies: actual full-environment kernel and scalar certificate theorem
% from the repository at 1bb9ebdc991090033b4d287411b66f056483f1d8.
% Labels use scp:. Requires standard theorem/proposition/corollary environments.

\section{Signed certificate polynomials with global noncancellation}
\label{scp:section}

Give the initial inert-newborn core genealogical depth zero. For each
future core $v$, let $\sigma_v$ be its birth time and let
$\operatorname{depth}(v)$ be its core-genealogical depth. Define
\[
 A_m^\infty(s)=\mathbb E_\infty
 \sum_{\operatorname{depth}(v)=m}e^{-s\sigma_v}.
\]
The existing full-environment characterization gives, for $s>0$,
\[
 r_\infty(s)=\sup_{m\ge1}
 \left(\frac{A_m^\infty(s)}{12(s+3)}\right)^{1/m},
 \qquad r_\infty(\mu)=1,
 \qquad A_m^\infty(s)\le s^{-m}.
\]
In particular $A_m^\infty(\mu)\le12(\mu+3)$, and
$s<\mu$ if and only if some finite retained-history sum exceeds
$12(s+3)$. We use precisely these established facts. The first-promotion
transform alone does not supply this characterization.

\subsection{Retained histories and their denominators}

Fix integers $m,R\ge1$ and $N\ge0$. Start from the actual inert
newborn, with one real core, one real leaf, and one inert parent stub.
The budget $R$ counts all actual core-clock rings, including null choices
of the stub. The budget $N$ counts all real vertex additions after this
two-vertex initial state. Let $L_{m,R,N}(s)$ be the expected sum of
$e^{-s\sigma_v}$ over depth-$m$ core births whose birth prefix has used
at most $R$ core rings and at most $N$ real additions. This is exactly
the positive finite-history lower sum obtained by retaining beta bursts
and core outcomes within these budgets. Rewards already earned before
the first discarded event are retained. A birth at the terminal retained
ring is counted when its population is still within budget.

At each retained ring, before leaf growth the skeleton has $k$ cores
and $\ell$ leaves, and the burst adds $b$ leaves. Its occupation entry is
\[
 R_{H,s}(L,M)=
 \left[\prod_{i:L_i>0}\binom{M_i-1}{L_i-1}\right]
 \frac{b!}{\prod_{u=0}^{b}(s+k+\ell+u)},
 \qquad b=|M|-\ell.
\]
Every displayed shift $k+\ell+u$ is the real population at a stage
of that burst, so it lies between $2$ and $N+2$.
For completeness, a zero-leaf state has occupation entry $1/(s+k)$;
this has the same stated form and range whenever such a state is used.

There is no additional uncancelled $k+s$ denominator from choosing the
next core. One can write the same discounted transition as the product
of the discounted ring mass $k/(k+s)$, the normalized burst law
$(k+s)R_{H,s}$, the uniform source-core probability $1/k$, and its
neighbour-choice probability. The first three factors multiply to
$R_{H,s}$. A specified core outcome has neighbour-choice probability
equal to an integer divided by its source degree $d_j$. In the one-stub
reachable state class,
\[
 2\le d_j\le N+2.
\]
Different source choices leading to the same retained state can be kept
separate for this argument and then summed.

Put
\[
 M_N=\operatorname{lcm}(1,2,\ldots,N+2),\qquad
 D_{R,N}(s)=M_N^R\prod_{j=2}^{N+2}(s+j)^R,
 \qquad d_{R,N}=R(N+1).
\]

\begin{lemma}[A positive coefficient clearer]\label{scp:clearer}
The polynomial
\[
 Q_{m,R,N}(s)=D_{R,N}(s)L_{m,R,N}(s)
\]
has nonnegative integer coefficients. Its degree is at most
$d_{R,N}$, and
\[
 Q_{m,R,N}(1)\le D_{R,N}(1).
\]
\end{lemma}
\begin{proof}
A retained reward prefix has some $r\le R$ core rings. Each burst
contains any factor $s+j$ at most once, so the product of all its
burst denominators contains that factor with multiplicity at most $R$.
Each neighbour degree divides $M_N$, so the product of the $r$ degree
denominators divides $M_N^R$. All remaining burst numerators, including
factorials and binomial coefficients, are nonnegative integers. After
multiplication by $D_{R,N}$, that prefix is therefore a polynomial with
nonnegative integer coefficients. There are finitely many retained
prefixes, even in the presence of null rings, because both budgets are
finite. Summing proves the coefficient assertion and degree bound.
Finally positivity and the potential-genealogy bound give
$L_{m,R,N}(1)\le A_m^\infty(1)\le1$.
\end{proof}

\subsection{An explicit omitted history proves strict nonvanishing}

Define the signed residual and its cleared polynomial by
\[
 F_{m,R,N}(s)=12(s+3)-L_{m,R,N}(s),\qquad
 P_{m,R,N}(s)=12(s+3)D_{R,N}(s)-Q_{m,R,N}(s).
\]
Also put
\[
 b_{m,R}(s)=
 \frac{[2(s+2)]^{-R}}{2^m(s+2)_m}.
\]

\begin{theorem}[Full-critical noncancellation]\label{scp:nonvanishing}
For every $m,R\ge1$ and $N\ge0$,
\[
 F_{m,R,N}(\mu)\ge b_{m,R}(\mu)>0,
 \qquad
 P_{m,R,N}(\mu)\ge D_{R,N}(\mu)b_{m,R}(\mu)>0.
\]
If $a\le\mu\le b$ is any positive rational enclosure, then the fully
rational lower bounds
\[
 F_{m,R,N}(\mu)\ge b_{m,R}(b),\qquad
 P_{m,R,N}(\mu)\ge D_{R,N}(a)b_{m,R}(b)
\]
hold. Moreover $P_{m,R,N}$ is an integer polynomial satisfying
\[
 \deg P_{m,R,N}\le R(N+1)+1,
 \qquad
 H(P_{m,R,N})\le49D_{R,N}(1)
 =49\left[M_N\frac{(N+3)!}{2}\right]^R.
\]
Here $H$ denotes the maximum absolute coefficient.
\end{theorem}
\begin{proof}
Consider the event prefix consisting of $R$ initial root rings choosing
the inert stub, followed by $m$ consecutive promotions along a single
descending core chain, with no other source-clock rings between these
specified events. Every initial null ring occurs in the unchanged
two-real-vertex state, so its discounted mass is
$1/[2(s+2)]$. Before the $(j+1)$st promotion, the real population is
$j+2$, and the terminal core has one parent and one leaf. Its specified
promotion has source rate one and leaf-choice probability $1/2$;
the discounted next-event mass is therefore $1/[2(s+j+2)]$.
Multiplication gives exactly $b_{m,R}(s)$.

The final birth has depth $m$, but occurs at core-ring number $R+m$,
so its reward is absent from $L_{m,R,N}$, regardless of $N$. All
other omitted rewards are nonnegative. Consequently
\[
 A_m^\infty(s)-L_{m,R,N}(s)\ge b_{m,R}(s),\qquad s>0.
\]
At $s=\mu$, combine this with
$A_m^\infty(\mu)\le12(\mu+3)$. Multiplication by the strictly positive
$D_{R,N}(\mu)$ proves the second assertion. The rational bounds use
that $b_{m,R}$ decreases and $D_{R,N}$ increases on the positive axis.

The integer and degree assertions follow from Lemma~\ref{scp:clearer}.
Since $D_{R,N}$ has nonnegative coefficients,
\[
 H(12(s+3)D_{R,N})\le48D_{R,N}(1),\qquad
 H(Q_{m,R,N})\le Q_{m,R,N}(1)\le D_{R,N}(1).
\]
The coefficientwise triangle inequality gives the stated height bound.
\end{proof}

This is noncancellation for a family of signed arithmetic expressions
attached to the full original-model critical relation. It does not use
an algebraicity hypothesis on $\mu$. The displayed height bound is a
coarse explicit bound, not an assertion that the chosen clearer or the
resulting polynomial is primitive or optimal.

\subsection{An increasing sequence of algebraic roots}

Every $L_{m,R,N}$ extends continuously to $s=0$, is nonincreasing on
$[0,\infty)$, and tends to zero at infinity. If it is not identically
zero it is strictly decreasing, because each retained reward is a
Laplace transform of a positive birth time. Thus $F_{m,R,N}$ is
strictly increasing and tends to infinity.

Define $\xi_{m,R,N}=0$ when $L_{m,R,N}(0)\le36$. Otherwise define
$\xi_{m,R,N}$ to be the unique positive solution of
\[
 L_{m,R,N}(s)=12(s+3),
\]
equivalently the unique positive zero of $P_{m,R,N}$.
The intermediate value theorem and the preceding monotonicity prove
existence and uniqueness in the second case.

\begin{theorem}[Exact algebraic-root representation]\label{scp:roots}
Every $\xi_{m,R,N}$ is a nonnegative algebraic number strictly below
$\mu$. For
\[
 \xi_h=\max_{1\le m\le h}\xi_{m,h,h}\qquad(h\ge1),
\]
one has
\[
 0\le\xi_h\le\xi_{h+1}<\mu,
 \qquad \boxed{\mu=\lim_{h\to\infty}\xi_h=\sup_{h\ge1}\xi_h.}
\]
Equivalently, $\xi_h$ is the largest nonnegative real root of the
explicit integer polynomial
\[
 \mathcal P_h(s)=s\prod_{m=1}^h P_{m,h,h}(s),
 \qquad \deg\mathcal P_h\le h^3+h^2+h+1.
\]
Each $\mathcal P_h(\mu)$ is strictly positive.
\end{theorem}
\begin{proof}
The coefficients of $P$ are integers, and its positive zero, when
present, lies strictly below $\mu$ by Theorem~\ref{scp:nonvanishing}.
The convention $\xi=0$ has the same strict upper bound.

Increasing either budget retains every previously retained reward and
possibly more. Hence, for fixed $m$, $L_{m,h,h}(s)$ increases with $h$
at every nonnegative $s$, and its threshold root is nondecreasing,
including the zero-root convention. The finite maximum over $m\le h$
is therefore nondecreasing.

Fix $0<s<\mu$. The established full scalar lower-certificate theorem
gives finite $m,R,N$ with $L_{m,R,N}(s)>12(s+3)$. For every
$h\ge\max\{m,R,N\}$, inclusion of retained histories gives
$L_{m,h,h}(s)>12(s+3)$, so $\xi_h>s$. Since every $\xi_h<\mu$,
this proves convergence. The product polynomial has no positive roots
other than the individual threshold roots, and the factor $s$ supplies
the root zero. Its degree and nonvanishing assertions follow by
multiplication.
\end{proof}

\begin{corollary}[Finite residuals for the exact pressure equation]
\label{scp:pressure-residual}
For $h\ge1$ and $s>0$, put
\[
 \rho_h(s)=\max_{1\le m\le h}
 \left(\frac{L_{m,h,h}(s)}{12(s+3)}\right)^{1/m},
 \qquad E_h(s)=1-\rho_h(s),
\]
using the nonnegative real root, including $0^{1/m}=0$.
Then
\[
 \rho_h(s)\uparrow r_\infty(s),\qquad
 E_h(s)\downarrow 1-r_\infty(s).
\]
The convergence is uniform on every compact subinterval of $(0,\infty)$,
without an asserted effective modulus in $h$. At the actual exponent,
\[
 E_h(\mu)>0\quad\hbox{for every finite }h,
 \qquad E_h(\mu)\longrightarrow0.
\]
In particular,
\[
 \boxed{\displaystyle
 \mu\text{ is the unique positive solution of }
 \lim_{h\to\infty}E_h(s)=0.}
\]
Each finite residual is obtained from rational functions by nonnegative
real radicals and a finite maximum, so it is algebraic at every positive
algebraic $s$.
\end{corollary}
\begin{proof}
Increasing $h$ enlarges both budgets and the range of generations in
the maximum, proving monotonicity. For each fixed $m$, retained-history
exhaustion and monotone convergence give
$L_{m,h,h}(s)\uparrow A_m^\infty(s)$: every birth prefix uses finitely
many real additions and core rings. Therefore
\[
 \lim_{h\to\infty}\rho_h(s)
 =\sup_{m\ge1}
 \left(\frac{A_m^\infty(s)}{12(s+3)}\right)^{1/m}
 =r_\infty(s).
\]
More explicitly, the finite maximum is bounded above by the displayed
supremum, while its limiting value is at least each fixed-$m$ term.
The omitted-history bound proves $L_{m,h,h}(s)<A_m^\infty(s)$.
Since the maximum is finite, this gives
$\rho_h(s)<r_\infty(s)$, in particular $E_h(\mu)>0$.
Each $\rho_h$ is continuous, and the limiting pressure is continuous;
the monotone convergence is thus uniform on compact intervals by
Dini's theorem. Finally the established strict pressure crossing
identifies the unique zero of $1-r_\infty(s)$ as $\mu$.
\end{proof}

This residual hierarchy connects the finite signed certificates directly
to the full critical equation. The finite stages are not asserted to
have small arithmetic height, and the radicals do not supply a new
small-form estimate. The exact limiting equation does not provide an
explicit cutoff-to-error bound for evaluating its left side.

The algebraic numbers $\xi_h$ can be isolated and compared by finite
integer-polynomial algorithms. No explicit bound on $\mu-\xi_h$ as a
function of $h$ follows from this proof. In particular, convergence of
these roots is not asserted to be rapid or computationally inexpensive.
The finite coefficient bound $L_{m,R,N}(0)\le R$ already shows that
no positive root is possible when $R\le36$. The following exact bound
gives a substantially stronger obstruction near the current enclosure.

\begin{proposition}[A necessary core-ring budget]\label{scp:ring-budget}
For $s>0$, define
\[
 p_i(s)=\prod_{j=1}^{i}\frac{j}{j+s},\qquad
 S_R(s)=\sum_{i=1}^{R}p_i(s).
\]
Every retained coefficient satisfies
\[
 L_{m,R,N}(s)\le S_R(s).
\]
The majorant has the exact evaluation
\[
 S_R(s)=
 \begin{cases}
 \displaystyle\frac{(R+1)p_R(s)-1}{1-s},&s\ne1,\\[6pt]
 H_{R+1}-1,&s=1,
 \end{cases}
\]
where $H_n=\sum_{j=1}^n1/j$. Any certificate root in
\[
 I=\left[\frac{113}{200},\frac{2832242711}{5000000000}\right]
\]
necessarily has $R\ge1222$.
\end{proposition}
\begin{proof}
Before the $i$th actual core-clock ring, the number $k_i$ of real
cores is at most $i$: the initial population has one core and every
preceding core ring adds at most one. Between these rings $k_i$ is
constant, regardless of leaf births. The next core-clock waiting time
can therefore be constructed as $E_i/k_i$, with fresh independent
mean-one exponentials $E_i$. If $\tau_i$ denotes the $i$th actual
core-ring time, then pathwise
\[
 \tau_i=\sum_{j=1}^{i}\frac{E_j}{k_j}
 \ge\sum_{j=1}^{i}\frac{E_j}{j}.
\]
Each retained generation-$m$ birth occurs at one of the first $R$
rings and contributes at most one discounted unit there. Ignoring its
generation and population restriction consequently gives
\[
 L_{m,R,N}(s)
 \le\sum_{i=1}^{R}\mathbb E e^{-s\tau_i}
 \le\sum_{i=1}^{R}\prod_{j=1}^{i}\frac{j}{j+s}=S_R(s).
\]
For $s\ne1$, the recurrence $(i+s)p_i=i p_{i-1}$ telescopes to
$(1-s)S_R=(R+1)p_R-1$. At $s=1$, $p_i=1/(i+1)$.

Put $a=113/200$. Exact rational evaluation of the same recurrence gives
the outward decimal enclosures
\[
 42.76694161<S_{1221}(a)<42.76694162,
 \qquad
 42.78297644<S_{1222}(a)<42.78297645.
\]
In particular $S_{1221}(a)<12(a+3)=2139/50=42.78$.
Since $S_R(s)$ increases with $R$ and decreases with $s$, if
$R\le1221$ and $s\in I$ then
\[
 L_{m,R,N}(s)\le S_{1221}(a)<12(a+3)\le12(s+3).
\]
Thus the defining root equation cannot hold there. The displayed
inequalities concern only a necessary budget; $R=1222$ is not asserted
to produce a successful root or certificate.
\end{proof}

\subsection{A quantified root gap and the remaining arithmetic task}

For any retained history weight, its logarithmic derivative in absolute
value is at most $d_{R,N}/(s+2)$. Summation therefore gives
\[
 0\le-L'_{m,R,N}(s)
 \le\frac{d_{R,N}}{s+2}L_{m,R,N}(s),\qquad s\ge0.
\]
If a positive root $\xi=\xi_{m,R,N}$ belongs to a rational enclosure
$[a,b]$ for $\mu$, then on $[\xi,\mu]$ the monotonicity of $F$ gives
$L(s)\le12(s+3)\le12(b+3)$. Consequently
\[
 12\le F'(s)\le
 K_{R,N}(a,b):=12+\frac{12d_{R,N}(b+3)}{a+2},
\]
and the mean value theorem yields the explicit nonzero gap
\[
 \boxed{\displaystyle
 \mu-\xi\ge\frac{b_{m,R}(b)}{K_{R,N}(a,b)}>0.}
\]
This lower bound is a separation statement; it is not the upper error
estimate needed for rapid approximation or a transcendence argument.

There is an exact diagnostic for the literal clearers chosen above.
When $R=N=h$ and $1\le m\le h$, cancellation of $(s+2)^h$ gives
\[
 D_{h,h}(s)b_{m,h}(s)
 =\frac{M_h^h}{2^{h+m}}
 \frac{\displaystyle\prod_{j=3}^{h+2}(s+j)^h}
      {\displaystyle\prod_{j=2}^{m+1}(s+j)}.
\]
For $s>0$, the product $\prod_{j=3}^{h+2}(s+j)$ is greater than
$\prod_{j=2}^{m+1}(s+j)$ and greater than one. Since $M_h\ge6$,
\[
 \boxed{P_{m,h,h}(\mu)>(3/2)^h,\qquad
 \mathcal P_h(\mu)>\mu(3/2)^{h^2}.}
\]
Thus these literal cleared values grow rather than form small arithmetic
linear forms. Dividing out content or positive common polynomial factors
may change their sizes; this diagnostic is not an impossibility theorem
for primitive normalization or other signed constructions.

The construction thus supplies explicit signed integer polynomials,
full-critical nonvanishing, a coarse degree/height bound, and a monotone
algebraic-root representation of the actual exponent. What remains
missing is a sufficiently small value relative to degree and height,
or another global arithmetic representation with such a bound. No
irrationality or transcendence of $\mu$ follows from the results here.
